% Options for packages loaded elsewhere
\PassOptionsToPackage{unicode}{hyperref}
\PassOptionsToPackage{hyphens}{url}
\PassOptionsToPackage{dvipsnames,svgnames,x11names}{xcolor}
\documentclass[
  11pt,
]{article}
\usepackage{xcolor}
\usepackage[margin=1in]{geometry}
\usepackage{amsmath,amssymb}
\setcounter{secnumdepth}{-\maxdimen} % remove section numbering
\usepackage{iftex}
\ifPDFTeX
  \usepackage[T1]{fontenc}
  \usepackage[utf8]{inputenc}
  \usepackage{textcomp} % provide euro and other symbols
\else % if luatex or xetex
  \usepackage{unicode-math} % this also loads fontspec
  \defaultfontfeatures{Scale=MatchLowercase}
  \defaultfontfeatures[\rmfamily]{Ligatures=TeX,Scale=1}
\fi
\usepackage{lmodern}
\ifPDFTeX\else
  % xetex/luatex font selection
\fi
% Use upquote if available, for straight quotes in verbatim environments
\IfFileExists{upquote.sty}{\usepackage{upquote}}{}
\IfFileExists{microtype.sty}{% use microtype if available
  \usepackage[]{microtype}
  \UseMicrotypeSet[protrusion]{basicmath} % disable protrusion for tt fonts
}{}
\makeatletter
\@ifundefined{KOMAClassName}{% if non-KOMA class
  \IfFileExists{parskip.sty}{%
    \usepackage{parskip}
  }{% else
    \setlength{\parindent}{0pt}
    \setlength{\parskip}{6pt plus 2pt minus 1pt}}
}{% if KOMA class
  \KOMAoptions{parskip=half}}
\makeatother
\setlength{\emergencystretch}{3em} % prevent overfull lines
\providecommand{\tightlist}{%
  \setlength{\itemsep}{0pt}\setlength{\parskip}{0pt}}
\usepackage{bookmark}
\IfFileExists{xurl.sty}{\usepackage{xurl}}{} % add URL line breaks if available
\urlstyle{same}
\hypersetup{
  pdftitle={No anomaly-free discrete symmetry splits the vector-like triplets and doublets of three 27s by one power of the breaking field},
  pdfauthor={Matthew Gelfman (ORCID 0009-0000-6402-2258); Melissa Behrends-Gelfman (ORCID 0009-0001-6281-1947)},
  colorlinks=true,
  linkcolor={Maroon},
  filecolor={Maroon},
  citecolor={Blue},
  urlcolor={Blue},
  pdfcreator={LaTeX via pandoc}}

\usepackage{amssymb}
\usepackage{newunicodechar}
\newunicodechar{±}{\ensuremath{\pm}}
\newunicodechar{²}{\ensuremath{^{2}}}
\newunicodechar{³}{\ensuremath{^{3}}}
\newunicodechar{·}{\ensuremath{\cdot}}
\newunicodechar{Δ}{\ensuremath{\Delta}}
\newunicodechar{Σ}{\ensuremath{\Sigma}}
\newunicodechar{Φ}{\ensuremath{\Phi}}
\newunicodechar{₂}{\ensuremath{_{2}}}
\newunicodechar{₃}{\ensuremath{_{3}}}
\newunicodechar{₆}{\ensuremath{_{6}}}
\newunicodechar{ℤ}{\ensuremath{\mathbb{Z}}}
\newunicodechar{∎}{\ensuremath{\blacksquare}}
\newunicodechar{−}{\ensuremath{-}}
\newunicodechar{∓}{\ensuremath{\mp}}
\newunicodechar{≈}{\ensuremath{\approx}}
\newunicodechar{≡}{\ensuremath{\equiv}}
\newunicodechar{≢}{\ensuremath{\not\equiv}}
\newunicodechar{≤}{\ensuremath{\le}}
\title{No anomaly-free discrete symmetry splits the vector-like triplets
and doublets of three 27s by one power of the breaking field}
\author{Matthew Gelfman (ORCID 0009-0000-6402-2258) \and Melissa
Behrends-Gelfman (ORCID 0009-0001-6281-1947)}
\date{25 November 2026}

\begin{document}
\maketitle

\subsection{Abstract}

Models with three generations of vector-like colour triplets D(3,1,\ensuremath{-}1/3)
and weak doublets L(1,2,\ensuremath{-}1/2), as in the 27 of E\ensuremath{_{6}} (Gürsey, Ramond and
Sikivie 1976), often need the doublets far lighter than the triplets. A
natural mechanism is a discrete symmetry \ensuremath{\mathbb{Z}}\textsubscript{N} under which
a scalar \ensuremath{\Phi} gives D its mass at dimension four, while L can take mass
only from \ensuremath{\Phi}\ensuremath{^{2}}/M, so that M\textsubscript{L} \ensuremath{\approx} M\textsubscript{D}\ensuremath{^{2}}/M.

We show that no flavour-blind \ensuremath{\mathbb{Z}}\textsubscript{N} does this without a
mixed \ensuremath{\mathbb{Z}}\textsubscript{N}--SU(3)\ensuremath{^{2}} anomaly. Requiring SM Yukawa couplings
with one Higgs doublet makes the quarks' contribution to that anomaly
vanish, which leaves 3\ensuremath{\cdot}q\textsubscript{\ensuremath{\Phi}} from the three D states.
Cancelling it forces \ensuremath{\Phi}\ensuremath{^{3}} to be neutral. Then \ensuremath{\Phi}\ensuremath{^{2}} carries the charge of \ensuremath{\Phi}*,
so L gets its mass at dimension four as well, and the hierarchy is lost.
The argument holds for every N and every lepton, Higgs or quark charge
assignment, and does not use the gravitational condition.

A brute-force scan of every flavour-blind assignment for N \ensuremath{\le} 24 confirms
this. With the weaker ``half-strength'' form of the gauge conditions
used in some of the literature, a single solution survives (\ensuremath{\mathbb{Z}}\ensuremath{_{6}}, with the
leptons at charge 3), and only if the gravitational condition is also
taken at half strength. We correct a miscounted gravitational anomaly
that circulated for this \ensuremath{\mathbb{Z}}\ensuremath{_{6}} case. Ways around the result need extra
fields, a Green--Schwarz mechanism, or explicit breaking. The last is
incompatible with using the symmetry for protection (Krauss and Wilczek
1989).

\subsection{1. Setup}

Take three generations with flavour-blind \ensuremath{\mathbb{Z}}\textsubscript{N} charges:
q\textsubscript{Q}, q\textsubscript{u}, q\textsubscript{d},
q\textsubscript{l}, q\textsubscript{e}, q\textsubscript{N} for the SM
fermions (right-handed fields counted by their own charges), h for the
Higgs doublet and q\textsubscript{\ensuremath{\Phi}} for a singlet scalar \ensuremath{\Phi}. Yukawa
invariance gives q\textsubscript{u} = q\textsubscript{Q} + h,
q\textsubscript{d} = q\textsubscript{Q} \ensuremath{-} h, q\textsubscript{e} =
q\textsubscript{l} \ensuremath{-} h and q\textsubscript{N} = q\textsubscript{l} + h,
with the sign conventions absorbed into h.

Let \ensuremath{\Delta}\textsubscript{D} = q(D\textsubscript{R}) \ensuremath{-} q(D\textsubscript{L})
and \ensuremath{\Delta}\textsubscript{L} = q(L\textsubscript{R}) \ensuremath{-} q(L\textsubscript{L}).
We require four things:

\begin{itemize}
\tightlist
\item
  \textbf{(a) D mass at dimension four:} \ensuremath{\Delta}\textsubscript{D} \ensuremath{\equiv}
  \ensuremath{\mp}q\textsubscript{\ensuremath{\Phi}} (mod N) and \ensuremath{\Delta}\textsubscript{D} \ensuremath{\not\equiv} 0.
\item
  \textbf{(b) L mass only at dimension five:} \ensuremath{\Delta}\textsubscript{L} \ensuremath{\not\equiv} 0 and
  \ensuremath{\Delta}\textsubscript{L} \ensuremath{\not\equiv} \ensuremath{\pm}q\textsubscript{\ensuremath{\Phi}}, but \ensuremath{\Delta}\textsubscript{L} \ensuremath{\equiv}
  \ensuremath{\mp}2q\textsubscript{\ensuremath{\Phi}}.
\item
  \textbf{(c) No baryon-number-violating operators with a lepton:} QQQL,
  u\textsuperscript{c} u\textsuperscript{c} d\textsuperscript{c}
  e\textsuperscript{c} and QQu\textsuperscript{c} e\textsuperscript{c}
  are all forbidden.
\item
  \textbf{(d) No mixed anomaly with SU(3)\ensuremath{^{2}}, SU(2)\ensuremath{^{2}} or gravity.}
\end{itemize}

\subsection{2. Anomaly conditions}

For left-handed Weyl fields, an instanton of SU(M) has 2T(r) zero modes
per field in representation r, so the 't Hooft vertex ('t Hooft 1976)
carries \ensuremath{\mathbb{Z}}\textsubscript{N} charge A\textsubscript{M} = \ensuremath{\Sigma}
q\textsubscript{i} \ensuremath{\cdot} 2T(r\textsubscript{i}). Anomaly freedom requires
A\textsubscript{M} \ensuremath{\equiv} 0 (mod N) (Ibáñez and Ross 1991; Banks and Dine
1992). For gravity, the corresponding condition is A\textsubscript{grav}
= \ensuremath{\Sigma} q\textsubscript{i} dim(r\textsubscript{i}) \ensuremath{\equiv} 0 modulo N, or modulo
N/2 for even N.

The SM quarks contribute 3(2q\textsubscript{Q} \ensuremath{-} q\textsubscript{u} \ensuremath{-}
q\textsubscript{d}) = 0 to A\ensuremath{_{3}} by Yukawa invariance. The D states
contribute 3(q(D\textsubscript{L}) \ensuremath{-} q(D\textsubscript{R})) =
\ensuremath{-}3\ensuremath{\Delta}\textsubscript{D} = \ensuremath{\pm}3q\textsubscript{\ensuremath{\Phi}}. Therefore

A\ensuremath{_{3}} = \ensuremath{\pm}3 q\textsubscript{\ensuremath{\Phi}}.

\subsection{3. Theorem}

\emph{No flavour-blind \ensuremath{\mathbb{Z}}\textsubscript{N} with three generations
satisfies (a), (b) and the strict condition A\ensuremath{_{3}} \ensuremath{\equiv} 0 (mod N).}

\emph{Proof.} A\ensuremath{_{3}} \ensuremath{\equiv} 0 requires 3q\textsubscript{\ensuremath{\Phi}} \ensuremath{\equiv} 0 (mod N), so
2q\textsubscript{\ensuremath{\Phi}} \ensuremath{\equiv} \ensuremath{-}q\textsubscript{\ensuremath{\Phi}}. Condition (b) needs
\ensuremath{\Delta}\textsubscript{L} \ensuremath{\equiv} \ensuremath{\mp}2q\textsubscript{\ensuremath{\Phi}} \ensuremath{\equiv} \ensuremath{\pm}q\textsubscript{\ensuremath{\Phi}}, which
contradicts \ensuremath{\Delta}\textsubscript{L} \ensuremath{\not\equiv} \ensuremath{\pm}q\textsubscript{\ensuremath{\Phi}}. \ensuremath{\blacksquare}

The proof uses neither (c) nor the SU(2)\ensuremath{^{2}} and gravitational conditions.
With N\textsubscript{g} generations it generalises: anomaly freedom
needs N\textsubscript{g} q\textsubscript{\ensuremath{\Phi}} \ensuremath{\equiv} 0, and the obstruction is
that 2q\textsubscript{\ensuremath{\Phi}} must not equal \ensuremath{\pm}q\textsubscript{\ensuremath{\Phi}} or 0. For
N\textsubscript{g} = 3 the only possibility is the one excluded above.

\subsection{4. Scan and the half-strength conventions}

An exhaustive scan over all (q\textsubscript{Q}, q\textsubscript{l}, h,
\ensuremath{\Delta}\textsubscript{D}, \ensuremath{\Delta}\textsubscript{L}) for N = 2 to 24 under (a) to (c)
gives:

\begin{itemize}
\tightlist
\item
  \textbf{Strict gauge conditions (mod N):} no solution, with the
  gravitational condition at either strength or omitted.
\item
  \textbf{Gauge conditions at half strength (mod N/2):} only \ensuremath{\mathbb{Z}}\ensuremath{_{6}}
  survives, and only if the gravitational condition is also taken mod
  N/2 or omitted. With the gravitational condition mod N, nothing
  survives.
\end{itemize}

For the \ensuremath{\mathbb{Z}}\ensuremath{_{6}} case (q\textsubscript{Q} = 0, q\textsubscript{l} = 3,
q\textsubscript{\ensuremath{\Phi}} = 1, \ensuremath{\Delta}\textsubscript{D} = \ensuremath{-}1, \ensuremath{\Delta}\textsubscript{L} =
\ensuremath{-}2), the sums are A\ensuremath{_{3}} = 3, A\ensuremath{_{2}} = 15 and A\textsubscript{grav} = 21. The
gravitational sum has been quoted as 30, which treats the lepton doublet
and e\textsuperscript{c} as cancelling. Counted field by field, the
lepton doublet, e\textsuperscript{c} and N\textsuperscript{c} give 2\ensuremath{\cdot}3 \ensuremath{-}
3 \ensuremath{-} 3 = 0 per generation, so the Standard Model part is 0 and the
vector-like part is 21. The Standard Model fields alone then pass every
half-strength condition, so the vector-like states are not needed for
anomaly freedom in that convention either.

Charge 3 for the leptons is lepton parity inside \ensuremath{\mathbb{Z}}\ensuremath{_{6}}. That is what
forbids the baryon-number-violating operators with a lepton; the ``81
flavour combinations'' sometimes quoted are one condition.

\subsection{5. Ways around}

\begin{itemize}
\tightlist
\item
  \textbf{Extra fields.} Extra fields charged under SU(3) can cancel A\ensuremath{_{3}},
  but they add matter.
\item
  \textbf{Green--Schwarz.} An axion-like field can shift to cancel the
  anomaly (Green and Schwarz 1984).
\item
  \textbf{Explicit breaking.} The symmetry can be broken explicitly, for
  example by Planck-suppressed operators. But a symmetry broken by
  gravity gives no protection, and the ``dimension-five gravitational
  bias'' sometimes used to remove \ensuremath{\mathbb{Z}}\textsubscript{N} domain walls cannot
  be combined with anomaly freedom.
\end{itemize}

Without one of these, the relation M\textsubscript{L} \ensuremath{\approx}
M\textsubscript{D}\ensuremath{^{2}}/M must be taken as an assumption about couplings,
not as a symmetry statement.

\subsection{Data and code}

\texttt{analysis/farside\_charge\_rule\_2026-11-25.py} and
\texttt{analysis/refinement\_checks\_2026-11-25.py} in the TRPG code
archive, doi:10.5281/zenodo.24329239. This result bears on one
assumption of the framework in doi:10.5281/zenodo.24330763.

\subsection{Acknowledgements}

Due to medical necessity, AI language-model assistants were used, under
the authors' direction, to draft and edit text from the authors' ideas,
to carry out and check calculations, and to search the literature; every
result was checked by rerunning the published code. The authors reviewed
all content and take full responsibility for it.

\subsection{References}

\begin{itemize}
\tightlist
\item
  Banks, T. and Dine, M. (1992). Note on discrete gauge anomalies.
  \emph{Phys. Rev.~D} 45, 1424.
\item
  Green, M. B. and Schwarz, J. H. (1984). Anomaly cancellation in
  supersymmetric D = 10 gauge theory and superstring theory. \emph{Phys.
  Lett. B} 149, 117.
\item
  Gürsey, F., Ramond, P. and Sikivie, P. (1976). A universal gauge
  theory model based on E6. \emph{Phys. Lett. B} 60, 177.
\item
  Ibáñez, L. E. and Ross, G. G. (1991). Discrete gauge symmetry
  anomalies. \emph{Phys. Lett. B} 260, 291.
\item
  Krauss, L. M. and Wilczek, F. (1989). Discrete gauge symmetry in
  continuum theories. \emph{Phys. Rev.~Lett.} 62, 1221.
\item
  't Hooft, G. (1976). Symmetry breaking through Bell--Jackiw anomalies.
  \emph{Phys. Rev.~Lett.} 37, 8.
\end{itemize}

\end{document}
